% TOP_PROBLEM: {"id": 30007153, "title": "Positive Einstein metric classification in dimension four", "category_id": 6, "category": {"id": 6, "name": "geometry", "display_name": "Geometry", "description": "Euclidean and non-Euclidean geometry, geometric structures.", "slug": "geometry", "order_index": 6, "created_at": "2026-07-31T15:26:25.671Z"}, "status": "open"}
% FIRST_POSTED: 2026-10-06
% ATTEMPT_STATUS: solved
% OPENAI_CLAIM: {"schema_version":1,"source_repo":"https://github.com/openai/math","source_commit":"adc7f1241b42e322a6451854ab7e4b4c146bf78a","release_date":"2026-10-06","match":"full","resolution":"proved","independently_reviewed":false,"second_pass":"agreed","families":[{"family":"348","title":"Nonnegative-curvature Einstein classification and an <i>L</i><sup>2</sup> topological gap","subject":"Differential geometry","manuscripts":[{"title":"Positively curved Einstein four-manifolds","pdf_path":"preprints/Positively-curved-Einstein-four-manifolds-September-23-2026/paper.pdf","readme_path":"preprints/Positively-curved-Einstein-four-manifolds-September-23-2026/README.md","date":"2026-09-23","theorem_ref":"preprints/Positively-curved-Einstein-four-manifolds-September-23-2026/build/sections/01-introduction.tex#thm:classification"}],"lean":{"doc_path":"lean/docs/348.md","comparators":["lean/ComparatorChallenges/EinsteinFour.json","lean/ComparatorChallenges/EinsteinFour.lean"],"declarations":[],"covers_main_theorem":true},"reasoning_trace":null}]}
\documentclass[11pt]{article}
\usepackage[margin=25mm]{geometry}
\usepackage{fontspec}
\setmainfont{Latin Modern Roman}
\usepackage{amsmath,amssymb,mathtools}
\usepackage{unicode-math}
\setmathfont{Latin Modern Math}
\usepackage{enumitem}
\usepackage{xurl,hyperref,bookmark}
\hypersetup{colorlinks=true,urlcolor=blue}
\setcounter{secnumdepth}{0}
\setlength{\parindent}{0pt}
\setlength{\parskip}{5pt}
\setlength{\emergencystretch}{3em}
\begin{document}
\section*{Positive Einstein metric classification in dimension four}

\subsection{OpenAI's claim}
\textit{Positive-Einstein four-manifold classification}: The paper claims that closed connected positively curved Einstein four-manifolds are, up to scaling, the round sphere, Fubini--Study projective plane, or round real projective space. The last possibility accommodates nonorientability.

\subsection{Outline of the argument}
\textit{Positive-Einstein four-manifold classification}: After normalizing the Einstein constant, coupled inequalities track the two Weyl blocks distinguished by the Hodge star. Weighted variance, Hessian, and moment comparisons reduce the analytic step to exact polynomial signs, forcing one block to vanish. Characteristic identities and a volume argument identify the oriented models, followed by an orientation-cover analysis of the possible quotient.

\subsection{Scope relative to this problem}
\textit{Positive-Einstein four-manifold classification}: The paper claims that closed connected positively curved Einstein four-manifolds are, up to scaling, the round sphere, Fubini--Study projective plane, or round real projective space. The last possibility accommodates nonorientability. The parameters, regularity assumptions and normalization in this question are retained in applying that assertion.

With the hypotheses and conventions above, the claimed result addresses the whole mathematical assertion on this page.

\subsection{Formal verification}
Family 348 describes Lean coverage for the strict-positive-curvature Einstein classification; zero-plane rigidity and the small-energy gap are separate.

\textit{Positive-Einstein four-manifold classification}: The manuscript links \texttt{EinsteinFour.lean}.

The documented selection includes this manuscript’s main assertion; that assertion retains the restrictions stated above.

This site has not run the Lean development. These links describe OpenAI’s selected formal statements and should be read with their stated scope.

\subsection{Status caveat}
OpenAI reports that an internal model produced the manuscripts. Their arguments have not been independently verified as proofs by this site. Any portion outside the documented formal scope may contain mistakes, and the catalogue label reports the announced mathematical claim.

\subsection{Sources}
\begin{itemize}
\item OpenAI, \emph{Positively curved Einstein four-manifolds}, 2026-09-23; citation key \texttt{OAI:Positively-curved-Einstein-four-manifolds-September-23-2026}. \href{https://github.com/openai/math/blob/adc7f1241b42e322a6451854ab7e4b4c146bf78a/preprints/Positively-curved-Einstein-four-manifolds-September-23-2026/paper.pdf}{PDF (GitHub)}; \href{https://github.com/openai/math/raw/adc7f1241b42e322a6451854ab7e4b4c146bf78a/preprints/Positively-curved-Einstein-four-manifolds-September-23-2026/paper.pdf}{PDF (direct)}; \href{https://github.com/openai/math/blob/adc7f1241b42e322a6451854ab7e4b4c146bf78a/preprints/Positively-curved-Einstein-four-manifolds-September-23-2026/README.md}{Manuscript page and citation}; \href{https://github.com/openai/math/blob/main/preprints/Positively-curved-Einstein-four-manifolds-September-23-2026/paper.pdf}{Latest version}.
\item Family 348: \href{https://github.com/openai/math/blob/adc7f1241b42e322a6451854ab7e4b4c146bf78a/lean/docs/348.md}{Lean scope}.
\item \href{https://github.com/openai/math/blob/adc7f1241b42e322a6451854ab7e4b4c146bf78a/lean/ComparatorChallenges/EinsteinFour.json}{Comparator EinsteinFour.json}.
\item \href{https://github.com/openai/math/blob/adc7f1241b42e322a6451854ab7e4b4c146bf78a/lean/ComparatorChallenges/EinsteinFour.lean}{Comparator EinsteinFour.lean}.
\item \href{https://github.com/openai/math/blob/adc7f1241b42e322a6451854ab7e4b4c146bf78a/CONTENTS.md}{OpenAI catalogue entry}.
\end{itemize}
\end{document}
